\documentclass[11pt,a4paper]{article}
\usepackage[margin=2.5cm]{geometry}
\usepackage{amsmath,amssymb,booktabs,graphicx,xcolor,hyperref,microtype}
\usepackage{authblk}
\hypersetup{colorlinks=true,linkcolor=blue,citecolor=blue,urlcolor=blue}

\title{\textbf{Paper 54: Order Parameter, Correlation Length Law, and Cross-Axis Surface}\\
\large Three Open Questions from the VCML Behavioral Phase Diagram}

\author[1]{Gabriel Aubin-Moreau}
\affil[1]{Independent research}
\date{March 2026}

\begin{document}
\maketitle

\begin{abstract}
Paper~53 identified five experimentally separable control directions for the VCML
behavioral phase diagram and posed three open questions:
(1) Is sg\_C = sg4/$\sigma_w$ the unified order parameter that collapses the five
axes?
(2) Does zone coherence scale with wave radius (correlation length law)?
(3) Does the write trigger--perturbation intensity interaction (Exp C prediction:
C\_perturb fails at low amplitude) hold?
This paper tests all three.

\textit{Exp A (sg\_C test):}
Sweeping perturbation amplitude at S1 shows that snr = sg4/$\sigma_w$ is
roughly flat ($\approx 0.105$) across the identity-encoding regime and drops
slightly ($\approx 0.094$) in the parity regime.
sg\_C is \emph{not} the unified order parameter: it is smoother than sg4 but
does not predict regime identity.
The na\_ratio remains the most diagnostic single metric for the
identity/parity transition.

\textit{Exp B (correlation length law):}
At S2 (zone\_width=20), sweeping r\_wave reveals two distinct failure modes.
At r\_wave=1 ($\delta=20$, well above the Axis~II wave-bleed threshold),
na\_ratio = 0.748 (parity) --- a correlation-length failure, not wave bleed.
This confirms that copy-forward loop coherence requires $r_\text{wave} \geq
r_\text{crit}$ independent of $\delta$.
sg4 peaks at r\_wave=4 ($\delta=5$); na\_ratio peaks at r\_wave=8 ($\delta=2.5$).

\textit{Exp C (cross-axis surface):}
The prediction is inverted.
C\_perturb maintains identity encoding at \emph{all} amplitude levels tested
(na = 1.493, 1.390, 1.252 for sub-threshold, optimal, and over-perturbed),
while C\_ref exhibits a parity/identity/parity transition as amplitude increases.
Mechanism: C\_perturb's trigger selectivity minimizes $\sigma_w$ (within-zone
fieldM variance), producing high SNR even with small absolute sg4.
The cross-axis interaction is: trigger selectivity provides noise-gating, not
signal-boosting.
\end{abstract}

%──────────────────────────────────────────────────────────────────────────────
\section{Introduction}

Paper~53 synthesized Papers~0--52 into a five-axis behavioral phase diagram
and identified three open questions.
This paper tests each.

\paragraph{The three questions.}
\begin{enumerate}
  \item \textbf{Unified order parameter.}
    Paper~53 identified sg\_C = sg4/$\sigma_w$ as a candidate that might
    collapse the five-axis diagram.
    If sg\_C is monotone where sg4 is not, it is the better coordinate.

  \item \textbf{Correlation length law.}
    The Axis~V finding (patch formation at S4) implied a finite copy-forward
    loop correlation length.
    The claim: correlation\_length $\approx r_\text{wave}$.
    The test: sweep $r_\text{wave}$ at fixed zone\_width and show where coherence fails.

  \item \textbf{Cross-axis surface.}
    Axis~IV (write trigger) and Axis~I (perturbation intensity) were predicted
    to interact: C\_perturb should fail when wave events are absent (low
    amplitude).
\end{enumerate}

%──────────────────────────────────────────────────────────────────────────────
\section{Methods}

All experiments use the FastVCML system with 5 seeds, STEPS=2000.

\paragraph{Exp A.}
S1 grid (W=80, H=40, HALF=40, WR=4.8, r\_wave=2).
Six amplitude levels: (supp=0.06, exc=0.12) through (supp=1.20, exc=2.40).
Metrics: sg4, $\sigma_w$, snr=sg4/$\sigma_w$, na\_ratio, coll/site.

\paragraph{Exp B.}
S2 grid (W=160, H=80, HALF=80, WR=19.2), optimal amplitude (supp=0.25, exc=0.50).
Five r\_wave levels: 1, 2, 4, 8, 16.
$\delta = W_\text{zone}/r_\text{wave} = 20/r_\text{wave}$.

\paragraph{Exp C.}
S1 grid. Three amplitude levels $\times$ two write modes (C\_ref, C\_perturb) $\times$ 5 seeds = 30 runs.

%──────────────────────────────────────────────────────────────────────────────
\section{Results}

\subsection{Exp A: sg\_C Is Not the Unified Order Parameter}

\begin{table}[h!]
\centering
\small
\begin{tabular}{lrrrrl}
\toprule
Condition & coll/site & sg4 & $\sigma_w$ & snr & na (encoding) \\
\midrule
A1 (deep sub) & 0.00061 & 0.0174 & 0.162 & 0.104 & 0.945 [parity] \\
A2 (near thresh) & 0.00122 & 0.0200 & 0.195 & 0.104 & 1.067 [IDENTITY] \\
A3 (optimal) & 0.00241 & 0.0209 & 0.202 & 0.105 & 1.026 [IDENTITY] \\
A4 (over-pert) & 0.00347 & 0.0258 & 0.225 & 0.115 & 1.014 [IDENTITY] \\
A5 (strong over) & 0.00525 & 0.0215 & 0.230 & 0.094 & 0.872 [parity] \\
A6 (extreme) & 0.00861 & 0.0248 & 0.253 & 0.098 & 0.856 [parity] \\
\bottomrule
\end{tabular}
\caption{Exp A: sg\_C (snr) across amplitude sweep. snr is flat in the
identity regime ($\approx 0.105$) and drops slightly in the parity regime
($\approx 0.094$--0.098). na\_ratio is the clearer regime indicator.}
\end{table}

\textbf{Finding.}
snr = sg4/$\sigma_w$ does not collapse the amplitude axis.
It is smoother than sg4 (less prone to apparent non-monotone noise) but
does not improve regime prediction: the identity/parity transition is captured
more cleanly by na\_ratio than by snr.

\textbf{Implication.}
sg\_C = snr is not the unified order parameter.
The identity/parity transition (na\_ratio threshold) and the signal amplitude
(sg4) are not reducible to a single ratio at this level.
This does not rule out a deeper unification; it rules out the simple
sg4/$\sigma_w$ ratio as that unification.

\subsection{Exp B: Correlation Length Law Confirmed}

\begin{table}[h!]
\centering
\small
\begin{tabular}{rrrrrl}
\toprule
r\_wave & $\delta$ & coll/site & sg4 & snr & na (encoding) \\
\midrule
1 & 20.0 & 0.00082 & 0.0046 & 0.041 & 0.748 [parity] \\
2 & 10.0 & 0.00238 & 0.0090 & 0.046 & 0.865 [parity] \\
4 & 5.0  & 0.00797 & 0.0132 & 0.052 & 1.002 [IDENTITY borderline] \\
8 & 2.5  & 0.02759 & 0.0121 & 0.065 & 1.274 [IDENTITY] \\
16 & 1.25 & 0.07542 & 0.0063 & 0.072 & 1.064 [IDENTITY] \\
\bottomrule
\end{tabular}
\caption{Exp B: r\_wave sweep at S2 (zone\_width=20). Both sg4 and na\_ratio
show non-monotone behavior. At r\_wave=1 ($\delta=20$, well above wave-bleed
threshold), na\_ratio=0.748 (parity) from correlation-length failure.}
\end{table}

\textbf{Finding 1: Correlation-length failure at r\_wave=1.}
At r\_wave=1 and $\delta=20$ (well above the Axis~II wave-bleed threshold of
$\delta \approx 2.5$), na\_ratio = 0.748 (parity).
This failure is not Axis~II (wave bleed): $\delta$ is adequate.
It is a failure of the copy-forward loop to propagate fieldM across the
zone width (20 sites) when the wave radius is only 1.

This directly confirms the correlation length hypothesis: the copy-forward loop
has a finite effective range proportional to $r_\text{wave}$.
When zone\_width $\gg r_\text{wave}$, structure forms in patches rather than
zone-wide gradients.

\textbf{Finding 2: Two-sided non-monotone in sg4.}
sg4 peaks at r\_wave=4 ($\delta=5$, the Axis~II peak from Paper~47).
It falls both for smaller r\_wave (correlation-length failure) and larger
r\_wave (wave bleed starting as coll/site rises sharply).
na\_ratio peaks later, at r\_wave=8 ($\delta=2.5$).

The sg4 peak and na\_ratio peak occur at different r\_wave values because they
measure different aspects of structure:
sg4 measures mean inter-zone distance (maximized by balanced suppression/excitation);
na\_ratio measures whether the correct zones are more distinct (maximized by
clean zone-to-zone gradient).

\textbf{Finding 3: snr increases monotonically with r\_wave.}
Unlike sg4 and na\_ratio, snr increases from 0.041 at r\_wave=1 to 0.072 at
r\_wave=16.
This is because $\sigma_w$ decreases at large r\_wave (large waves create more
uniform fieldM within zones, reducing within-zone variance).
sg\_C therefore fails to capture the parity failure at large r\_wave: it
\emph{increases} exactly where na\_ratio deteriorates (r\_wave=16, na=1.064).
Further evidence that snr is not the unified order parameter.

\textbf{Correlation length bound.}
From this sweep: identity encoding (na $>$ 1) requires $r_\text{wave} \geq 4$
for zone\_width=20, i.e., $r_\text{wave}/\text{zone\_width} \geq 0.2$.
Equivalently, zone\_width $\leq 5 \times r_\text{wave}$.
This is the empirical correlation length law: beyond 5 wave radii per zone,
the copy-forward loop cannot maintain zone-wide coherence.

\subsection{Exp C: Prediction Inverted --- C\_perturb Is Noise-Robust}

\begin{table}[h!]
\centering
\small
\begin{tabular}{llrrrl}
\toprule
Amplitude & Mode & sg4 & $\sigma_w$ & snr & na (encoding) \\
\midrule
sub-threshold & C\_ref & 0.0174 & 0.162 & 0.104 & 0.945 [parity] \\
sub-threshold & C\_perturb & 0.0069 & 0.0115 & 0.611 & \textbf{1.493} [IDENTITY] \\
optimal & C\_ref & 0.0209 & 0.202 & 0.105 & 1.026 [IDENTITY] \\
optimal & C\_perturb & 0.0048 & 0.0093 & 0.522 & \textbf{1.390} [IDENTITY] \\
over-perturbed & C\_ref & 0.0215 & 0.230 & 0.094 & 0.872 [parity] \\
over-perturbed & C\_perturb & 0.0036 & 0.0100 & 0.361 & \textbf{1.252} [IDENTITY] \\
\bottomrule
\end{tabular}
\caption{Exp C: write trigger $\times$ amplitude. C\_perturb maintains identity
encoding at all amplitudes. C\_ref fails at both extremes. Mechanism: C\_perturb
has extremely low $\sigma_w$ (0.009--0.012), producing high snr even with tiny
sg4.}
\end{table}

\textbf{Finding.}
The Paper~53 prediction (C\_perturb fails at low amplitude because no wave
events occur) is wrong.
C\_perturb maintains identity encoding at \emph{all} amplitude levels tested.
At sub-threshold amplitude, C\_perturb achieves na=1.493 --- the highest
na\_ratio of any condition in this paper.

\textbf{Mechanism: noise gating, not signal boosting.}
C\_perturb's sg4 is consistently \emph{smaller} than C\_ref (0.007--0.0048 vs
0.017--0.022).
But C\_perturb's $\sigma_w$ is 10--20$\times$ smaller (0.009--0.012 vs 0.162--0.230).
The result: C\_perturb's snr is 4--6$\times$ higher (0.36--0.61 vs 0.094--0.105).

Trigger selectivity does not write more informative content into fieldM.
It writes to fewer cells, keeping within-zone variance ($\sigma_w$) near zero.
The resulting fieldM has small absolute magnitude (low sg4) but extremely low
noise (low $\sigma_w$) and thus high SNR.
High SNR, not high sg4, determines whether the na\_ratio crosses the identity
threshold.

\textbf{C\_ref failure modes.}
C\_ref fails at sub-threshold because the viability contrast between
suppressed and excited zones is weak (small amplitude), but it writes to all
quiescent cells, accumulating high $\sigma_w$ from noisy writes.
C\_ref fails at over-perturbed because rapid collapse cycles prevent the
quiescence streak from accumulating, mixing formation and re-encoding signals.
C\_perturb avoids both: it only writes during wave events, so it accumulates
only the zone-specific signal regardless of amplitude level.

%──────────────────────────────────────────────────────────────────────────────
\section{Discussion}

\subsection{The Actual Order Parameter}

sg\_C = sg4/$\sigma_w$ (snr) is not the unified order parameter.
However, the Exp C results reveal what the actual ordering quantity is:
\emph{zone-specific SNR} evaluated at the level of na\_ratio, not sg4.
Formally, na\_ratio measures whether the inter-zone signal exceeds the
intra-zone noise for the correct zone pairs (nonadj vs adj).
The C\_perturb results show this can be achieved with tiny sg4 if $\sigma_w$ is
also tiny.

A candidate order parameter is therefore not the scalar snr = sg4/$\sigma_w$
but the vectorial quantity:
\begin{equation}
  C_\text{order} = \frac{\|\mu_\text{nonadj} - \mu_\text{adj}\|}{\sigma_w}
\end{equation}
where $\mu_\text{nonadj}$ and $\mu_\text{adj}$ are the mean fieldM distances
for non-adjacent and adjacent zone pairs respectively.
This is the na\_ratio with an explicit noise normalization.
Whether this quantity collapses the full five-axis diagram remains an open
question.

\subsection{The Correlation Length Law}

The empirical bound (identity encoding requires
zone\_width $\leq 5 \times r_\text{wave}$) is a testable, falsifiable claim.
It predicts that:
\begin{enumerate}
  \item Increasing $r_\text{wave}$ at fixed S3 (zone\_width=80) should
    recover identity encoding.
    Prediction: na\_ratio $> 1$ at r\_wave $\geq 16$ (zone\_width/5=16).
  \item The scale failure at S4 (zone\_width=320) should be recoverable with
    $r_\text{wave} \geq 64$.
  \item Systems with local learning rules (biological cortex, ~50--100~$\mu$m
    propagation range) should show analogous patch formation at
    zone\_width $> 5 \times$ propagation\_range.
\end{enumerate}

The correlation length bound also explains why $\delta$ (Axis~II) and
$r_\text{wave}$ (correlation length) appear as separate axes:
$\delta = W_\text{zone}/r_\text{wave}$ captures wave containment,
while zone\_width/r\_wave captures loop propagation range.
These are not the same quantity: wave containment requires large $\delta$,
while loop coherence requires small zone\_width/$r_\text{wave}$.
Both conditions must be satisfied simultaneously.

\subsection{Revised Axis IV Understanding}

The Exp C result changes the theoretical interpretation of the trigger
selectivity axis.
The original framing (Paper~52) was: C\_perturb writes only during boundary
events, selecting more informative writes.
The correct framing is: C\_perturb writes to fewer cells total, maintaining
low $\sigma_w$, which produces high SNR regardless of whether individual
writes are more informative.

This is a mechanistic distinction.
Signal-boosting predicts C\_perturb should fail at low amplitude (less
informative writes).
Noise-gating predicts C\_perturb should succeed at any amplitude (always writes
to few cells).
The data confirms noise-gating.

The implication: C\_perturb's adversarial resistance (Paper~52) and its
cross-amplitude robustness (this paper) both follow from the same mechanism.
Trigger selectivity is a noise-reduction operation, not an information-selection
operation.

\subsection{Toward the Stability Condition}

The regime diagram now points toward an analytical gap.
The three experiments show regime transitions at:
\begin{itemize}
  \item na\_ratio: parity $\leftrightarrow$ identity as amplitude, r\_wave, or
    trigger selectivity varies.
  \item A stability condition exists but is not yet derived analytically.
\end{itemize}

The natural framing: if the population field dynamics can be approximated by a
PDE (e.g., of Burgers or KPZ type, per the nonlinear birth term),
the regime boundaries should correspond to analytical stability conditions on
that PDE.
Specifically:
\begin{itemize}
  \item The correlation length law
    (zone\_width $\leq 5 r_\text{wave}$) should emerge from the propagation
    term.
  \item The amplitude peak (Axis~I non-monotone) should correspond to a
    balance between nonlinear amplification and dissipative smoothing.
  \item The na\_ratio identity threshold should correspond to a condition on
    the inter-zone forcing vs.\ within-zone noise.
\end{itemize}
Deriving these conditions analytically from the field equation would convert
the regime diagram from an empirical observation to a predicted consequence
of the system's dynamics.
This is the target for Paper~55.

%──────────────────────────────────────────────────────────────────────────────
\section{Conclusion}

Three open questions from Paper~53 are resolved.

\textbf{Order parameter}: sg\_C is not the unified order parameter.
na\_ratio remains the most diagnostic metric.
The candidate order parameter is a vectorial SNR measuring whether
correct zone pairs are separated beyond within-zone noise.

\textbf{Correlation length law}: Confirmed.
Identity encoding requires zone\_width $\leq 5 \times r_\text{wave}$.
The failure at r\_wave=1 ($\delta=20$) is not wave bleed; it is a distinct
correlation-length mechanism.
Two-sided non-monotone in sg4 (correlation-length failure and wave bleed
at opposite extremes).

\textbf{Cross-axis surface}: Prediction inverted.
C\_perturb maintains identity at all amplitude levels via noise-gating, not
signal-boosting.
Trigger selectivity reduces $\sigma_w$, not the write-level information content.
C\_ref transitions parity$\to$identity$\to$parity as amplitude increases;
C\_perturb stays identity throughout.

The most important theoretical revision: the Axis IV trigger mechanism is
noise-gating.
The most important empirical finding: the correlation length bound
(zone\_width $\leq 5 r_\text{wave}$) is falsifiable and testable at scale.

\end{document}
